\section{Results and Evaluation}
\label{sec:results_2}

Models are evaluated on the ATIS dataset using Intent Classification Accuracy and CoNLL-style Slot F1. Table \ref{tab:2a_results} reports the results obtained by incrementally modifying the baseline architecture, while Table \ref{tab:2b_results} summarizes the performance of the fine-tuned pre-trained models.

\begin{table}[H]
\centering
\small
\caption{Incremental optimization results for Part 2.A (From Scratch).}
\label{tab:2a_results}
\resizebox{\columnwidth}{!}{%
\begin{tabular}{lcc}
\hline
\textbf{Incremental Modification Variant} & \textbf{Intent Acc.} & \textbf{Slot F1} \\ \hline
Provided Baseline ($d_{\text{model}}\!=\!20, L\!=\!1, H\!=\!1, \text{FF}_{\text{dim}}\!=\!20$) & 0.8847 & 0.8093 \\
+ Capacity Scale Up ($d_{\text{model}}\!=\!64, \text{FF}_{\text{dim}}\!=\!128$) & 0.9250 & 0.8834 \\
+ Full Capacity Scaling ($d_{\text{model}}\!=\!128, \text{FF}_{\text{dim}}\!=\!192$) & 0.9317 & 0.8839 \\
+ Multi-Head Expansion ($H\!=\!2, d_{\text{model}}\!=\!128$) & 0.9183 & 0.8721 \\
+ Deep Layer Stacking Expansion ($L\!=\!4, d_{\text{model}}\!=\!64$) & 0.9037 & 0.8674 \\
+ Depth Scaling Layer Setup ($L\!=\!2, d_{\text{model}}\!=\!64$) & 0.9272 & 0.8714 \\
+ Deep Multi-Head Scale Up ($L\!=\!4, H\!=\!2, d_{\text{model}}\!=\!64$) & 0.9306 & 0.8884 \\
+ Deep Regularized Layout ($\text{Dropout}_{\text{out}}\!=\!0.1, L\!=\!4, H\!=\!2, d_{\text{model}}\!=\!128$) & \textbf{0.9395} & \textbf{0.8906} \\ \hline
\end{tabular}
}
\end{table}

The baseline model achieves an Intent Accuracy of 0.8847 and a Slot F1 score of 0.8093. Increasing the model dimensions provides the largest improvement, with both tasks benefiting from the additional representational capacity. Scaling $d_{\text{model}}$ from 20 to 64 and later to 128 consistently improves performance, suggesting that model width is one of the main factors affecting results on ATIS.

Increasing the number of attention heads or the number of layers in isolation does not always produce further gains. Some deeper or more complex configurations show slightly lower scores, likely because of the relatively small size of the dataset. A more balanced setup combining increased capacity, additional layers, and multiple attention heads performs better than these isolated modifications.

The best model trained from scratch combines larger dimensions, deeper layers, multiple attention heads, and an output dropout of 0.1. This configuration reaches an Intent Accuracy of 0.9395 and a Slot F1 score of 0.8906, representing the strongest performance among the custom architectures.

\begin{table}[H]
\centering
\small
\caption{Multi-task transfer fine-tuning results for Part 2.B.}
\label{tab:2b_results}
\begin{tabular}{lcc}
\hline
\textbf{Pre-trained Backbone Model} & \textbf{Intent Acc.} & \textbf{Slot F1} \\ \hline
GPT2-small Initial Trial Configuration & 0.9474 & 0.8822 \\
GPT2 Fully Fine-Tuned Setup & 0.9586 & 0.8949 \\
BERT-base-uncased Architecture (Run 1) & 0.9653 & 0.9307 \\
BERT-base-uncased Architecture (Run 2) & \textbf{0.9765} & \textbf{0.9332} \\ \hline
\end{tabular}
\end{table}

Fine-tuning pre-trained transformer models produces substantially better results than any configuration trained from scratch. Even the initial \texttt{GPT2-small} setup exceeds the performance of the best custom architectures after only a few epochs of training. Fully fine-tuning GPT2 further improves performance, reaching an Intent Accuracy of 0.9586 and a Slot F1 score of 0.8949.

Among all evaluated models, \texttt{BERT-base-uncased} achieves the highest overall results. The best run obtains an Intent Accuracy of 0.9765 and a Slot F1 score of 0.9332. These improvements are likely related to BERT's bidirectional attention mechanism, which allows token representations to incorporate information from both left and right context. This is particularly useful for slot filling, where predictions often depend on surrounding words. In contrast, GPT2 relies on a causal attention mask and therefore only has access to preceding context when building token representations.

% --- FIRST PLOT CONTAINER (INTENT & SLOT METRICS ACROSS BOTH SYSTEMS) ---
\begin{figure}[h]
\centering
\begin{subfigure}{\columnwidth}
  \centering
  \includesvg[width=0.88\columnwidth]{images/intent_acc_test_2a_2b}
  \caption{Multi-Task Intent Accuracy Evaluation Profile (Combined Test Set)}
  \label{fig:intent_acc_2a_2b}
\end{subfigure} \\
\vspace{3mm}
\begin{subfigure}{\columnwidth}
  \centering
  \includesvg[width=0.88\columnwidth]{images/slot_f1_test_2a_2b}
  \caption{Multi-Task Slot F1 Sequence Tagging Evaluation Profile (Combined Test Set)}
  \label{fig:slot_f1_2a_2b}
\end{subfigure}
\caption{Comparative performance metrics between custom models built from scratch and pre-trained backbone transformers, illustrating intent classification trajectories (a) and sequence-tagging token F1 accuracy gains (b).}
\label{fig:scratch_vs_pretrained_metrics}
\end{figure}

Figure \ref{fig:scratch_vs_pretrained_metrics} summarizes the performance differences between the models trained from scratch and the fine-tuned transformer backbones. The pre-trained models consistently achieve higher scores on both tasks, highlighting the benefit of transfer learning and large-scale pre-training.

The advantage is particularly visible for slot filling in Figure \ref{fig:slot_f1_2a_2b}, where \texttt{BERT-base-uncased} establishes a clear performance gap over the custom architectures. While the scratch-trained models improve through architectural scaling, they do not reach the same level of performance as the pre-trained alternatives.

% --- SECOND PLOT CONTAINER (CONVERGENCE & REGULARIZATION COEFFICIENTS) ---
\begin{figure}[h]
\centering
\begin{subfigure}{\columnwidth}
  \centering
  \includesvg[width=0.88\columnwidth]{images/loss_train_2a_2b}
  \caption{Cross-Entropy Convergence Profile (Combined Training Track)}
  \label{fig:loss_train_2a_2b}
\end{subfigure} \\
\vspace{3mm}
\begin{subfigure}{\columnwidth}
  \centering
  \includesvg[width=0.88\columnwidth]{images/loss_dev_2a_2b}
  \caption{Cross-Entropy Generalization Profile (Combined Validation Track)}
  \label{fig:loss_dev_2a_2b}
\end{subfigure}
\caption{Joint optimization loss trajectories across all evaluated configurations on both training (a) and validation (b) splits.}
\label{fig:cross_entropy_comparisons}
\end{figure}

Figure \ref{fig:cross_entropy_comparisons} compares the training and validation loss curves of all evaluated systems. Models trained from scratch require more optimization steps before reaching stable performance, whereas the fine-tuned transformer models converge much more rapidly.

This behaviour reflects the advantage of transfer learning: the pre-trained models start from useful linguistic representations and therefore require fewer updates to adapt to the ATIS tasks. As a result, both GPT2 and BERT achieve lower validation losses and more stable convergence than the custom architectures.